\section{Justification}
\label{sec:justification}

The generator standardizes the setup work that otherwise gets repeated manually for every report.

\begin{table}[H]
  \centering
  \begin{tabular}{@{}lll@{}}
    \toprule
    Component & Decision & Motivation \\
    \midrule
    Document class & report & Chapter-oriented thesis layout \\
    Bibliography & biblatex + biber & Stable citation workflow \\
    Build tool & latexmk & Correct multi-pass compilation \\
    Source layout & chapter folders & Maintainable long-form writing \\
    \bottomrule
  \end{tabular}
  \caption{Key document decisions inherited from the reference project. \textit{Source: Compiled by the author.}}
  \label{tab:template-decisions}
\end{table}

\begin{definition}
  A reusable LaTeX scaffold is a project skeleton that preserves formatting and build conventions while minimizing setup cost.
\end{definition}
\loesource{Compiled by the author.}

\begin{figure}[H]
  \centering
  \begin{tikzpicture}[node distance=2.3cm, >=Latex]
    \node[draw, rounded corners, fill=blue!5, minimum width=3.1cm, minimum height=1cm] (template) {Builder};
    \node[draw, rounded corners, fill=red!5, right=of template, minimum width=3.3cm, minimum height=1cm] (generator) {Scaffold Script};
    \node[draw, rounded corners, fill=green!5, right=of generator, minimum width=3.1cm, minimum height=1cm] (project) {New Report};
    \draw[->, thick] (template) -- (generator);
    \draw[->, thick] (generator) -- (project);
  \end{tikzpicture}
  \caption{Flow from reusable builder to generated LaTeX project. \textit{Source: Compiled by the author.}}
  \label{fig:scaffold-flow}
\end{figure}
